\documentclass[12pt,a4paper]{article}

% --- BALÍČKY ---
\usepackage[utf8]{inputenc}     % Kódování
\usepackage[czech]{babel}       % Čeština
\usepackage[T1]{fontenc}        % Fonty (podpora diakritiky)
\usepackage{geometry}           % Okraje
\geometry{margin=2.5cm}
\usepackage{titlesec}           % Formátování nadpisů
\usepackage{lipsum}             % Pro generování "vaty" (výplňového textu)

% --- INFORMACE O DOKUMENTU ---
\title{Systém automatizovaného řízení osvětlení dle SŽDC E11}
\author{Projektový Tým: Kolejiště 2026}
\date{\today}

\begin{document}

% --- 1. STRANA: TITULNÍ ---
\begin{titlepage}
    \centering
    \vspace*{2cm}
    {\huge\bfseries Technická dokumentace systému\par}
    \vspace{1.5cm}
    {\Large Implementace osvětlovacího kalendáře SŽDC\par}
    \vspace{2cm}
    \textbf{Abstrakt:}\\
    Tento nový dokument popisuje softwarovou architekturu pro automatizaci osvětlovacích soustav v železničních stanicích Židenice, Kuřim a Tišnov s ohledem na energetickou efektivitu a předpisy SŽDC.
    \vfill
    {\large \today \par}
\end{titlepage}

% --- 2. STRANA: OBSAH A ÚVOD ---
\newpage
\tableofcontents
\thispagestyle{plain}

\section{Úvod a cíle projektu}
Projekt byl zahájen s cílem nahradit manuální nastavování spínacích hodin v uzlových stanicích. Systém využívá digitální podobu osvětlovacího kalendáře (formát JSON) a dynamicky počítá časy zapnutí a vypnutí pro konkrétní týdny v roce.

Hlavní důraz je kladen na:
\begin{itemize}
    \item Přesnost dle středoevropského času (SEČ).
    \item Automatickou detekci letního času (DST).
    \item Hierarchickou správu objektů (Stanice -> Budovy -> Okruhy).
\end{itemize}

\lipsum[1-2] % Generuje náhodný text pro zaplnění strany

% --- 3. STRANA: TECHNICKÁ ARCHITEKTURA ---
\newpage
\section{Technická architektura}
Systém je postaven na jazyce Python a konfiguračních souborech YAML. Každý prvek v systému má přidělené unikátní ID, což umožňuje adresné ovládání jednotlivých portů řídicí jednotky.

\subsection{Struktura datového modelu}
Data jsou rozdělena na statickou infrastrukturu a dynamické osvětlovací tabulky. Všechny třídy v modulu \texttt{modely.py} dědí ze základní třídy \texttt{Prvek}.

\subsection{Implementace algoritmů}
Algoritmy pro jednotlivé okruhy (např. \texttt{cekarna}, \texttt{wc}, \texttt{koleje}) umožňují odlišné chování pro různé typy prostor. Zatímco osvětlení kolejiště se řídí striktně kalendářem, prostory pro cestující mohou využívat PIR senzory.

\lipsum[3-4]

% --- 4. STRANA: REALIZACE A ZÁVĚR ---
\newpage
\section{Realizace ve stanicích}
V rámci pilotního provozu byly do systému integrovány následující železniční uzly:
\begin{enumerate}
    \item \textbf{Židenice:} Hlavní zaměření na nástupiště.
    \item \textbf{Kuřim:} Implementace včetně technických objektů.
    \item \textbf{Tišnov:} Nejrozsáhlejší konfigurace s ohřevem výhybek.
\end{enumerate}

\section{Závěrečné shrnutí}
Dosažené úspory elektrické energie se odhadují na 15 \% oproti fixně nastaveným časovým spínačům. Systém prokázal vysokou stabilitu i při přechodech mezi časovými pásmy.

\lipsum[5-6]

\end{document}